%% kapa: exam=1 round=31 year=2020 period=1 form=A subject=economics subject_ko=경제학원론 seq=38 number=78
%% source: Q-Net 감정평가사 기출문제 (cst00309) artlSeq=5207219 file=extracted/1008/5207219/A형/A형/2020년도 제31회 감정평가사 1차 시험_1교시_A형.pdf
%% conversion: kordoc 4.24.51 -> tools/gen1.py (auto); review=proofread-sonnet; flags=none
\begin{question}{78}
경기변동이론에 관한 설명으로 옳은 것은?
\begin{choices}
  \item 실물경기변동(real business cycle)이론에서 가계는 기간별로 최적의 소비 선택을 한다.
  \item 실물경기변동이론은 가격의 경직성을 전제한다.
  \item 실물경기변동이론은 화폐의 중립성을 가정하지 않는다.
  \item 가격의 비동조성(staggering pricing)이론은 새고전학파(New Classical) 경기변동이론에 속한다.
  \item 새케인즈학파(New Keynesian)는 공급충격이 경기변동의 원인이라고 주장한다.
\end{choices}
\end{question}
\answer{1}
